\chapter{Related Work}
\label{sec:related}

Mechanistic interpretability work on arithmetic has focused on \textit{how} transformers solve 
these tasks internally. \citet{nanda2023progress} studied a one-layer transformer trained on 
modular addition and found that grokking corresponds to the emergence of a Fourier-based 
algorithm, given that the model encodes operands as superpositions of trigonometric features and composes 
them across layers. Related work on pretrained language models finds similarly structured 
numerical representations, suggesting frequency-domain arithmetic features are not limited to 
small synthetic settings~\citep{zhou2024fourier}.

\citet{zhong2023clock} showed that modular addition admits two distinct solutions --- a 
Fourier-based \textit{clock} algorithm and a lookup-table \textit{pizza} algorithm --- with which emerges 
depending on training configuration. This demonstrates that identical supervision can produce 
structurally different circuits. For standard integer addition, \citet{quirke2024understanding} 
identified a carry-propagation circuit across attention and MLP layers resembling long addition, 
and \citet{musat2024clustering} showed that such algorithmic structure emerges through sharp 
phase transitions rather than gradual accumulation.

On the representation domain, \citet{park2023linear} formalise the linear representation 
hypothesis, under which concepts correspond to directions in the residual stream. 
\citet{zou2025representationengineeringtopdownapproach} implement this through 
representation engineering, extracting concept directions from contrastive prompt sets and 
applying them as targeted interventions. \citet{engels2024notall} show that not all features 
are one-dimensional, motivating the use of broader localisation methods beyond single directions.

On knowledge editing, \citet{meng2023locatingeditingfactualassociations} locate factual 
associations within specific MLP layers and edit them directly. Soft prompting 
\citep{panickssery2024steeringllama2contrastive} offers a parameter-free alternative as learned 
prefix vectors shift model behaviour without modifying weights, making it appropriate for concepts resistant to layer-level localisation.

The present work asks whether mechanistic structure comparable to that found in small or 
proprietary models is recoverable in \texttt{Qwen3-4B-Instruct}. 